% language=us runpath=texruns:manuals/ontarget

\startcomponent ontarget-mathsizes

\environment ontarget-style

\startchapter[title={The three math sizes}]

What follow is a short explanation about how a \TEX\ math engine deals with
sizes. When typesetting a formula there are three different sizes assumed, as
shown in the next example:

\startlinecorrection
\scale[scale=3000]{\im{x^{x^x}}}
\stoplinecorrection

If we take each \im{x} and scale it to the same height we get the following. It
is clear that the smallest size runs wider and looks a bit bolder which makes it
more readable at the smaller size.

\startlinecorrection
\dontleavehmode
\scale[height=2cm]{\im{\textstyle         x}}
\scale[height=2cm]{\im{\scriptstyle       x}}
\scale[height=2cm]{\im{\scriptscriptstyle x}}
\stoplinecorrection

So why is this? The \im {x} is a math italic character and in traditional \TEX\
that comes from what is called a math italic font. However, when \TEX\ processes
the list of intermediate math nodes, that \im {x} is specified as two numbers: a
family and an index. A family is a collection of three references to fonts: they
can be the same or different but whatever one chooses, they have to be scaled to
what is expected. The initial \TEX\ setup for instance makes a family from \type
{cmmi10}, \type {cmmi7} and \type {cmmi5}. Each font has a slightly different
design and when we typeset at 10 points no scaling takes place.

That all changed a bit when \OPENTYPE\ math fonts came around. In fact, there was
only one such font to start with so that basically set the standard: Cambria from
\MICROSOFT. Instead of making three variants, the fonts contains upto three
alternatives (although there is no real limit) for a specific shape. These
variants are \type {ssty=1} and \type {ssty=2} where \type {ssty} refers to the
alternative substitution font feature.

Fonts like Latin Modern also follow this model. So, in the case of an \im {x}, we
have three options to choose from: the normal shape, or one of the \type {ssty}
replacements. Because that font is mostly compatible with Computer Modern the \im
{x} is like \type {cmmi10} and the smaller sizes like \type {cmmi} and \type
{cmmi5} respectively. It is important to keep in mind that we are talking of one
font! That means that in principle we only consume one family. Where with eight
bit fonts we need quite some families, for instance for all these symbols, with
\OPENTYPE\ fonts in a \UNICODE\ universe we can stick to one. Of course you can
still use the more traditional model of multiple families, but just a copy of
this one font will do then. In practice most symbols scale perfectly well to a
smaller size so no special design size variants are provided and that makes an \type
{ssty} mapping table relatively small: only for characters that benefit from it
get an alternate.

It should also be noted that the alternates are not scaled down in their design,
so actually the script \im {x} is in the font as if it were from \type {cmmi7}
scaled up to 10 points. The font specifies two scaling factors but a macro
package can decide to scale differently. \footnote{Many font parameters in an
\OPENTYPE\ math fonts can be considered to have recommended values, and sometimes
it makes sense to change them.} All this means that in \TEX\ a family has three
times the same font, the original (text), and two scaled down copies with \type
{ssty=1} replacements (script and scriptscript). Although in practice an \type
{ssty} feature will provide two alternates, there is no rule that mandates that.
This introduces a dilemma:

\startnarrower
What are we supposed to do when there is a \type {ssty=1} but no \type {ssty=2}:
do we use a scaled down text for scriptscript or do we scale down the script
variant? Because all this is very much design related one can debate this.
\stopnarrower

Actually this is a general question. In for instance the Latin Modern Math font
there is an \type {aalt} feature where some glyphs have two and others three
variants. Should choosing variant three when there are two available use the
first or the last. The fact that in \type {ssty} there is some order, from small
to even smaller, other alternate features provide just that: an alternative
shape. It is good to keep in mind that \OPENTYPE\ provides mechanisms, just that,
and these can be uses any way a designer wants. This is why we see the ligature
mechanism being used for non|-|ligatures and ligatures being constructed by
contextual replacements (that still can be single characters).

After dealing with math fonts for decades one gets a feeling of what is needed by
users. And after implementing math font mechanism for eight bit fonts as well as
\OPENTYPE\ fonts, one also gets an idea how much the ideal situation differs from
reality. You can make a subsystem with plenty degrees of freedom but in the end
when little gets used it only adds overhead. And math actually comes with some.
In a traditional engine one has to deal with a limited number of families, which
results in switching mechanisms (for instance to full range bold, not only
alphabets) that also redefine math characters and reassign families. Because even
a simple \type {$x$} already can trigger quite some initialization one has to be
careful in coding this. For instance what happens with \type {{\bf $x$}} or \type
{{\small $x$}} is not just a simple family switch. In for instance section titles
switching style and size happens at the same time. I won't go into details how we
dealt with all that in \CONTEXT, lets stick to: we just have to make sure that
the overhead is acceptable.

When working on \LUAMETATEX\ it was decided to introduce a more dynamic scaling
model. First for regular text, later for math. There are several reasons for
this. For text if can reduce the overhead of loading fonts at different sizes. It
saves memory as well as load time, which involves scaling and passing to \TEX,
something that is noticeable for large fonts used for non Latin scripts.

This is also true for math but there we can go further. If instead of defining a
font three times (per family) we pass the engine the information where a script
and from that a scriptscript variant is located in the font, then we can delay
scaling till such a glyph is needed. Every character has a \type {smaller} field
and because we also support right to left math typesetting there's even a \type
{mirror} field. It is up to the engine to deal with that. In principle we could
now even turn family id's into font id's, but then we also need to handle the way
families are used in math character specifications. For that reason mid 2023 this
hasn't yet happened. We also don't want to abandon this nice \TEX\ family
concept, if only because we always need to support both models.

So, to summarize this:

\startitemize[packed]
    \startitem
        The traditional approach to sizes is to have three fonts with possibly
        different designs for shapes. A family consists of three different fonts.
    \stopitem
    \startitem
        The \OPENTYPE\ approach is to have optional \type {ssty} alternates for
        different shapes at smaller sizes. A family can use the same font scaled
        differently.
    \stopitem
    \startitem
        The \LUAMETATEX\ approach is to have one font for al three with size
        variants as part of the glyph definition. The engine replaces and scales
        on demand.
    \stopitem
\stopitemize

It will be clear that in the case of \LUAMETATEX, when scaling is delayed, the
engine has to do more work, because not only glyphs scale but also related font
parameters. And because we can also apply dynamic scaling of text to math, that
means multiple times scaling every time a dimension is needed. Not only that, we
also scale in two dimensions so for every math parameter we need to take that
into account. It meant a rather fundamental overhaul of the code, on top of the
plenty extra features that the \LUAMETATEX\ math engine provides. Of course we
could not divert too much from the principe approach to math typesetting, which
is why this all happened stepwise and why it is an option. The fact that we have
dynamic scaling also means that we can have smaller than scriptscript sizes
without much effort, but that is besides the point we want to make here. Anyway,
when a math font is defined one can set a \typ {compactmath} field as well as a
\typ {textscale}, \typ {scriptscale} and \typ {scriptscriptscale} when defining
the font (via the \LUA\ interface)

So far what we do is not that much different from what fonts have been doing
since the beginning of \TEX. Even the potential lack of \type {ssty=2} dilemma is
not unique: one can have design sizes eight bit fonts with holes too and it is
actually harder to deal with that. However, there are other potential issues.
Just like \type {ssty} is an alternate feature there can be more. And alternates
are a great way out of endless discussions about preferred shapes. Some math
fonts therefore have many alternates in the \type {ssXX} and \type {cvXX} feature
namespace. Of course there is no agreement about this so there is our second
dilemma:

\startnarrower
    How do we deal with additional alternates, given that the distribution of
    alternative shapes is arbitrary as is their grouping.
\stopnarrower

Although one way to control this is using \UNICODE\ modifiers, that is not
something that \TEX ies like to input. It still demands some agreement about
modifier related to original shapes and it doesn't cover the plenty unspecified
ones. And how portable is this?

If we assume that a math user wants consistency (which is quite an assumption)
one solution is to specify additional substitutions when we load the font. But then
we have a third dilemma:

\startnarrower
    Which substitution comes first: the chosen style alternatives or the size
    related alternates?
\stopnarrower

One can argue that the order that features are defines in fonts also determine
this accumulated substitutions. That makes a lot of sense but when you look at
some specification of how scripts are supposed to be dealt with in \OPENTYPE\ you
can find somewhat vague remarks about the order in which specific features are to
be applied. However, in practice one can best follow the order in the font. Say that
we want a different lowercase delta. There are two options:

\startitemize[packed,n]
\startitem
    The stylistic variant comes before the sizing one.
\stopitem
\startitem
    The sizing variant comes before the stylistic one.
\stopitem
\stopitemize

Now, if fonts are well designed, when there is a different design for each size,
there should also be one for the variant. Or the reverse, when there are
different variants, there should be different sizes. In practice that means that
there can be six different deltas. The replacement tables can be made such that
they both take care of mapping. Think of these glyphs names being related:

\starttyping
delta        delta.ss03
delta.ssty1  delta.ss03.ssty1
delta.ssty2  delta.ss03.ssty2
\stoptyping

So what if the two deltas are missing we first replace the sizes? We basically get this:

\starttyping
delta        delta.ss03
delta.ssty1  delta.ss03 (scaled)
delta.ssty2  delta.ss03 (scaled)
\stoptyping

or this:

\starttyping
delta        delta.ss03
delta.ssty1  delta (scaled)
delta.ssty2  delta (scaled)
\stoptyping

It depends bit om the machinery and the order of features in the font. It makes
no sense to point fingers to fonts here, simply because what would be shown here
could be different after a font gets updated.

In \CONTEXT\ we have two choices of dealing with this. We can immediately apply
features, or we can delay them which permits a mixture of different deltas. This
is actually needed for calligraphic versus script alphabets as one document can
have both while \UNICODE\ math combines them, which then means that we need to
revert to a feature applied. \footnote {In \CONTEXT\ we have a variant on this
where we have a dedicated additional alphabets to this which also permits for
taking calligraphic or script from other fonts, because not all fonts have both.}

Even if we delay there is an issue: in the three fonts per family approach we
already have the \type {ssty} applied. So then we replace the delta based on the
size replacement. So, we can end up with:

\starttyping
text         : delta              -> delta.ss03
script       : delta.ssty1 scaled -> delta.ssty1 scaled
scriptscript : delta.ssty2 scaled -> delta.ssty1 scaled
\stoptyping

However, when we delay scaling and use one font, as described for \LUAMETATEX, we
haven't yet applied \type {ssty}, so instead of pre|-|scaled replacements:

\starttyping
text         : delta
script       : delta scaled -> delta.ssty1 scaled
scriptscript : delta scaled -> delta.ssty2 scaled
\stoptyping

We can first replace and then scale.

\starttyping
text         : delta    delta.ss03
script       : delta -> delta.ss03 -> scaled
scriptscript : delta -> delta.ss03 -> scaled
\stoptyping

In the non|-|replaced delta case we would have had:

\starttyping
text         : delta
script       : delta.ssty1 -> scaled
scriptscript : delta.ssty2 -> scaled
\stoptyping

The difference is that because we haven't yet done the scaling and \type {ssty}
replacement we check against the original \UNICODE\ and when there is a lacking
\type {delta.ss03.ssty1} at that size, we just scale down the text size \type
{delta.ss03}. One can argue that we then might violate the lookup because maybe
there was on purpose no alternate but how to distinguish design from error in
fonts?

Of course when replacement is delayed one can also again consult the \type {ssty}
table and sort of fix the font runtime but there is a catch there: you end up
with a scaled glyph (scriptscript) in a differently scaled (text) font so you get
different dimensions and although the \CONTEXT\ backend can handle that it is
confusing and not pretty. One has to use the virtual font mechanism and in fact
the glyph we're dealing with can already be virtual (tweaked).

There are more cases where confusion can kick in. Examples are primes where can
have different shapes and positioning in text compared to the script sizes, and
accents, where text and replacements have either or not zero widths. To some
extend it is no real surprise because one needs a math engine that can expose it
when making the font. And after a while the font can't change any more because
it's being and there are many \TEX\ users who expect the same output after years.
I can imagine that a text editor that basically assumes one specific math font
being used (or expects other math fonts to be set up the same way) can provide a
list of magic key combinations that choose some shape and|/|or provide some
graphical user interface. I remember presentations at Bacho\TeX\ where was
demonstrated how a desktop publishing application deals with alternates and every
version got different or additional interfaces for that. In such a case a feature
and interface depend on each other, something that is not the case in a \TEX\
ecosystem: there we just access whatever we want. Of course there we also need
some interface but there are at least some traditions in place, like naming
characters or hiding complex and inconsistent choices in macros.

% emptyset in stix
% plank constant in ebgaramond

\stopchapter

\stopcomponent
